\chapter{Quiz Answers}
\label{app:answers}

\section*{Chapter 1 --- Spring Boot}
\begin{enumerate}
  \item The env var wins: env vars outrank the config-server document,
  which outranks the packaged YAML. Right for containers because the
  image must be immutable across environments; the runtime injects
  what differs.
  \item \code{@Configuration}, \code{@ComponentScan},
  \code{@EnableAutoConfiguration}. The last: auto-configuration is
  conditional on the classpath, so removing a dependency silently
  removes the beans it triggered.
  \item It resolves to the literal string \code{placeholder}. The
  service can therefore \emph{start} without Google credentials ---
  OAuth fails at click time instead of boot time.
  \item Exposure makes the endpoint exist; Spring Security still
  authorizes each request. If \code{/actuator/**} is not permitted,
  probes get 401, which the kubelet counts as failure.
\end{enumerate}

\section*{Chapter 2 --- Config Server}
\begin{enumerate}
  \item Env var $>$ config-server $>$ packaged YAML.
  \item Native serves files from the classpath directory inside the
  jar; git adds audit, review, revert, and config changes without
  rebuilding the config-server.
  \item \code{optional:} lets a service boot when the config-server is
  down --- but then it boots on localhost defaults and fails later,
  far from the cause.
  \item In its own \code{application.yml}, packaged in its jar. The
  bootstrapping principle: the source of configuration cannot
  configure itself; every such chain needs a self-contained root.
\end{enumerate}

\section*{Chapter 3 --- Eureka}
\begin{enumerate}
  \item course-service registers (POST) at startup; the gateway's
  discovery client refreshes its registry cache (default up to
  30\,s); routing then load-balances from cache. Worst case for a new
  instance to receive traffic: registration interval + cache refresh,
  tens of seconds.
  \item Compose makes service \emph{names} resolvable, and hostnames
  match them; in Kubernetes the pod hostname is the pod name, which
  no one can resolve --- so the IP must be registered instead.
  \item Availability of registry data over accuracy. Helps: a switch
  reboot silences all heartbeats but instances still run. Hurts: mass
  real crashes leave dead instances registered.
  \item The Service's endpoint list replaces the registry; readiness
  probes replace heartbeats; kube-proxy/DNS replace client-side
  lookup; the ReplicaSet replaces self-registration of new instances.
\end{enumerate}

\section*{Chapter 4 --- API Gateway}
\begin{enumerate}
  \item Browser $\rightarrow$ ts.local (hosts file) $\rightarrow$
  Traefik (Ingress: \code{/api} $\rightarrow$ api-gateway Service)
  $\rightarrow$ gateway pod (predicate \code{/api/auth/**},
  StripPrefix=2, JWT not required for \code{/api/auth}) $\rightarrow$
  Eureka-resolved auth-service pod, arriving as \code{/login}
  $\rightarrow$ controller.
  \item Google validates the redirect URI byte-for-byte; it was
  registered with the full \code{/login/oauth2/...} path, so it must
  arrive unmodified. API routes strip because controllers are written
  prefix-less.
  \item Every request with a token signed by auth-service fails the
  gateway's validation --- universal 401s on protected routes while
  login itself appears to work.
  \item A proxy is I/O-bound with high concurrency --- event loops
  excel. course-service does blocking JDBC work; reactive would
  complicate it without benefit until proven otherwise.
\end{enumerate}

\section*{Chapter 5 --- Security}
\begin{enumerate}
  \item CSRF rides an ambient credential (the session cookie sent
  automatically). A JWT in an Authorization header is attached
  explicitly by code; a forged cross-site request cannot add it.
  \item Fifteen minutes (its expiry). Nothing --- a signature cannot be
  revoked. \code{/refresh}: the refresh token is checked against the
  database, where revocation lives.
  \item Invariant: every network path to a service passes the
  gateway's JWT filter. Breaks: an Ingress routing straight to a
  service (the deleted \code{/courses} route), or any in-cluster
  caller forging headers on the flat pod network.
  \item The client-secret never leaves servers; the authorization
  code transits the browser. The code is single-use, short-lived, and
  worthless without the client-secret at exchange time.
  \item HS256: one shared secret; anyone who can validate can also
  sign. RS256: private key signs, public key validates --- so an IdP
  can let thousands of services validate without any of them being
  able to mint tokens.
\end{enumerate}

\section*{Chapter 6 --- Services and Common}
\begin{enumerate}
  \item (a) HTTP --- the caller needs the data now. (b) Event --- a
  fact others react to; grading must not wait on mail. (c) Neither:
  validation is local cryptography with the shared secret; no call.
  \item auth-service only \emph{produces} --- sends fail per-call,
  boot succeeds. notification-service only \emph{consumes}; without a
  broker it is pure idle cost, so it ships with Kafka.
  \item Belongs: event records --- both sides must agree on shape.
  Never: business logic/entities --- shared behavior couples services'
  release cycles and hides the boundary.
  \item Consumers deploy before producers, and changes are
  additive-only (new optional fields). Then old consumers ignore new
  fields and new consumers tolerate old events.
\end{enumerate}

\section*{Chapter 7 --- Postgres, JPA, Flyway}
\begin{enumerate}
  \item Two developers' schemas silently diverging (update applies
  per-database, no history); and destructive changes being impossible
  --- update never drops/renames, so entity refactors quietly desync
  from the schema.
  \item Each history row stores a checksum of its script; an edited
  script mismatches and Flyway refuses to start --- the history is
  the audit trail, and edits would falsify it.
  \item Release 1: add new column, write both (or trigger). Backfill
  migration. Release 2: read new column, stop writing old.
  Release 3: drop old column. Each step compatible with the version
  before it.
  \item For: fault and blast-radius isolation --- one server's crash,
  upgrade, or misconfiguration cannot touch the other's data. Against:
  double memory and operational surface for a homelab. Managed clouds
  usually take credential-isolation on one instance because servers
  are the billing unit.
\end{enumerate}

\section*{Chapter 8 --- MongoDB}
\begin{enumerate}
  \item Dictionary entries: self-contained, read/written as a unit,
  naturally nested, no cross-entry consistency. Enrollments: reference
  two other aggregates and require consistency between them --- joins
  and transactions are the point.
  \item Into the application code (and optional collection
  validators). Disciplines: backward-readable document classes (new
  fields optional), and explicit data-fix scripts for shape changes.
  \item Postgres refuses connections to a nonexistent database;
  Mongo creates databases on first write, so a wrong name yields
  empty results. Loud --- failures you notice at deploy beat failures
  you notice at audit.
  \item ClusterIP-only + no Ingress means only in-cluster callers can
  reach it; the compose parity also keeps dev friction zero.
  Invalidated the moment it is exposed beyond the cluster or shares
  the cluster with untrusted workloads.
\end{enumerate}

\section*{Chapter 9 --- MinIO / S3}
\begin{enumerate}
  \item The service signs (microseconds) instead of streaming
  (seconds); storage serves bytes directly. New failure mode: the
  URL embeds an endpoint and expiry --- wrong endpoint or clock skew
  breaks downloads in the browser only, invisible to the service.
  \item The signing service reaches MinIO as \code{minio:9000}
  (cluster DNS); the browser cannot resolve that name. The signature
  covers the host, so rewriting after signing invalidates it ---
  hence signing against a second, public endpoint.
  \item Expiry (15 minutes) bounds the damage window; the URL grants
  one object, not the bucket; and it cannot be re-signed without the
  keys.
  \item Changes: endpoint(s), credentials (ideally to an IAM role),
  durability guarantees. Unchanged: every SDK call, bucket/key/presign
  logic --- the API is the standard.
\end{enumerate}

\section*{Chapter 10 --- Kafka Core}
\begin{enumerate}
  \item Key by userId --- per-key order is per-partition order.
  Partitions $\geq$ max consumers you ever want in the group;
  parallelism is capped by partition count.
  \item Groups are independent: each keeps its own offsets per
  partition, so both groups traverse the full log; consumption is
  non-destructive reading plus a per-group cursor.
  \item \code{LISTENERS} = sockets bound; \code{ADVERTISED\_LISTENERS}
  = addresses \emph{returned to clients} in metadata. Bootstrap uses
  the address you supplied; all subsequent connections use the
  advertised ones --- so the failure begins after a successful hello.
  \item Retries and post-processing crashes both re-deliver;
  the broker cannot know if you acted. Notification: dedupe by event
  id (or accept a rare duplicate mail --- idempotence by tolerance).
\end{enumerate}

\section*{Chapter 11 --- Spring Kafka}
\begin{enumerate}
  \item \code{send()} buffers client-side and returns a future; no
  broker needed until delivery. The outage surfaces per-send: futures
  fail after timeouts, and buffered events are lost if the process
  dies.
  \item With no committed offsets, \code{earliest} starts from the
  oldest retained record --- history is processed. Once the group
  commits offsets, the setting is never consulted again.
  \item Without: deserialization throws inside poll, offset never
  advances, the partition stalls on one record forever. With: the
  error is materialized as a typed failure, the error handler skips
  (or retries then routes) --- in full production, to
  \code{<topic>.DLT}.
  \item JSON deserialization instantiates classes named in headers; an
  unrestricted trust list lets a message make the consumer construct
  arbitrary classpath types --- a gadget-injection vector.
  \item When ``DB write + event publish'' must not diverge. It makes
  the state change and the \emph{recording of the intent to publish}
  atomic --- one local transaction; a relay makes the publish
  eventually happen.
\end{enumerate}

\section*{Chapter 12 --- Docker}
\begin{enumerate}
  \item Namespaces (isolated views) and cgroups (resource limits).
  No guest kernel: millisecond starts and no per-instance OS cost ---
  but also: the kernel is shared (weaker isolation than a VM) and the
  JVM sees host resources unless told otherwise.
  \item Inside a container \code{localhost} is the container itself.
  Compose fixes it with service-name DNS on the bridge network;
  Kubernetes with Service DNS --- both via env-var overrides of the
  localhost defaults.
  \item Each container has its own network namespace, so both bind
  ``their'' 5432. Conflicts return at the shared host: published
  host ports (hence 5433/5434) --- and with \code{--network host},
  where there is no namespace separation at all.
  \item Plain \code{depends\_on} waits for \emph{started};
  \code{service\_healthy} waits for the healthcheck to pass ---
  here, config-server's curl of its own \code{/actuator/health},
  exposed by its management config.
\end{enumerate}

\section*{Chapter 13 --- Jib and the Registry}
\begin{enumerate}
  \item No daemon in CI, volatility-based layering (small re-pushes),
  non-root/OpenShift-ready images. The security one: no daemon means
  no docker socket in CI --- the socket is root on the host.
  \item The image name doubles as pull address for the cluster. They
  agree here because Jenkins shares the host's network namespace; if
  Jenkins were bridge-networked, its \code{localhost:5000} would be
  itself and pushes would fail while pulls worked.
  \item An import per node per deploy, minutes of serialization, and
  no immutable tag history; the pipeline's \code{set image} to a
  SHA-tagged name has nothing to point at --- the registry is what
  makes tags addressable.
  \item Same tag, new content: with \code{IfNotPresent} the node's
  cached image wins silently. Fixes: \code{imagePullPolicy: Always}
  (the manifests' choice), or deploy by immutable tags/digests
  (the pipeline's choice) --- this project does both.
\end{enumerate}

\section*{Chapter 14 --- Cluster Anatomy}
\begin{enumerate}
  \item You write desired state; controllers reconcile actual toward
  it forever. A second apply produces zero diff, so zero action ---
  idempotence by architecture.
  \item Ordering is a procedural concept; convergence replaces it:
  everything starts, dependents crash and retry with backoff until
  dependencies exist. CrashLoop during convergence is the mechanism,
  not a bug.
  \item kubectl sends a REST delete to the API server; state store
  updates; the ReplicaSet controller's loop sees replicas 0 $<$ 1 and
  creates a pod object; the scheduler assigns it; the kubelet on the
  node starts the container.
  \item Apply merges desired state; it never deletes what is absent
  from input. \code{kubectl delete ingress course-service -n ts}, or
  an apply with \code{--prune} (or a GitOps controller) would.
\end{enumerate}

\section*{Chapter 15 --- Workloads}
\begin{enumerate}
  \item Each template change births a new ReplicaSet while old ones
  are retained at scale 0 --- rolling update and one-command rollback
  fall out of keeping the counter separate from the versioned
  template.
  \item Deployment: N interchangeable replicas, any can die, order
  irrelevant. StatefulSet: named replicas with per-identity storage
  and ordered replacement. Eureka: Deployment (registry rebuilt by
  re-registration). Redis cache: Deployment if truly cache.
  Kafka broker: StatefulSet. Gateway: Deployment.
  \item New pods must pass readiness before old pods terminate ---
  Chapter~17's probes gate Chapter~15's rollout.
  \item A Deployment update runs old and new pods concurrently; two
  postgres processes opening one data directory corrupt it. The
  StatefulSet replaces one-at-a-time and re-binds the claim to the
  successor, never two at once.
\end{enumerate}

\section*{Chapter 16 --- Config and Secrets}
\begin{enumerate}
  \item RBAC on who may read Secret objects; lifecycle separation
  from git; encryption at rest (optional). This project: RBAC yes
  (Jenkins' Role omits secret-read... it grants create/patch but the
  practical control is namespace scoping), git-separation yes (script),
  encryption at rest no.
  \item Server image reads \code{POSTGRES\_*}; Spring reads
  \code{SPRING\_DATASOURCE\_*}. Split Secrets mean each pod imports
  only names it understands --- no foreign variables that mask typos
  or leak scope.
  \item Update: \code{JWT\_SECRET=new ./create-secrets.sh}; then
  immediately \code{kubectl -n ts rollout restart deployment/auth-service
  deployment/api-gateway}. Half-done: signer and validator hold
  different keys --- every login 401s, or worse, only after the next
  unrelated restart.
  \item The config-server's \code{configurations/*.yml} tree, mounted
  as files --- retiring the config-server itself.
\end{enumerate}

\section*{Chapter 17 --- Probes and Resources}
\begin{enumerate}
  \item Startup failing past budget: restart. Readiness failing:
  removed from endpoints, no restart --- safe to be aggressive.
  Liveness failing: restart --- the one to be conservative with.
  \item \code{/v3/api-docs} answers only after the full context
  (Flyway included) is up and is permitAll'd; a bound port proves
  only Tomcat exists. \code{/actuator/health} is 401 where security
  does not permit it (course/auth) and honest where it does
  (dictionary, gateway).
  \item The kernel's OOM killer terminated the cgroup for exceeding
  its memory limit --- no JVM involvement, so no exception. Check:
  the limit vs.\ \code{MaxRAMPercentage} arithmetic, and whether
  load or a leak grew the heap.
  \item Scheduler reads requests (placement); kernel enforces limits.
  Memory is incompressible --- exceeding it kills someone, so bound
  it predictably; CPU is compressible --- limits would throttle
  startup exactly when the JVM needs burst.
\end{enumerate}

\section*{Chapter 18 --- Networking}
\begin{enumerate}
  \item DNS resolves the Service name to a stable virtual IP; the
  endpoint controller keeps the Ready-pod list current; kube-proxy
  (or k3s' equivalent) rewrites virtual-IP traffic to a chosen pod
  IP; readiness decides membership.
  \item Only Ready pods are endpoints, so a rolling update shifts
  traffic exactly when new pods pass probes and old ones terminate
  --- zero-downtime is the probe-endpoint interaction.
  \item Selector/label mismatch --- the Service matches zero pods.
  Nothing errors because a Service with no endpoints is a legal
  state; the symptom is only visible in \code{get endpoints}.
  \item \code{ts.local} is not real DNS; the hosts file pins it to
  the tailnet IP. Replacements: a real domain with actual DNS +
  cert-manager TLS, or a Cloudflare Tunnel publishing the hostname.
\end{enumerate}

\section*{Chapter 19 --- Storage}
\begin{enumerate}
  \item Pods mount PVCs; PVs (and StorageClasses) are cluster-scoped;
  the indirection lets the same claim get local-path here and EBS on
  EKS --- manifests unchanged.
  \item Postgres runs on one node at a time --- RWO suffices. Shared
  uploads push people to RWX; object storage is usually the right
  answer instead.
  \item StatefulSet deletion spares PVCs so data survives
  orchestration mistakes; PVC deletion (reclaim \code{Delete})
  removes PV and data --- claims, not workloads, are where data
  lives.
  \item Persistence: survives restarts (yes). Durability: survives
  hardware failure (no --- one SSD). Backup: survives your own
  mistakes (no --- absent, the honest gap).
\end{enumerate}

\section*{Chapter 20 --- Kustomize}
\begin{enumerate}
  \item A definitive set (nothing forgotten, one command) and
  cross-cutting edits (namespace, labels) stamped consistently.
  \item Namespace: cluster-scoped --- Jenkins' namespaced Role cannot
  apply it. Secrets: values must not exist in git at all --- shape in
  script, values in the operator's environment.
  \item Kustomize: deploying your own app as readable YAML. Helm:
  distributing parameterized software to other people's clusters.
  Arriving as charts: Strimzi and kube-prometheus-stack.
  \item Nothing --- apply does not prune; mongodb keeps running,
  now unmanaged. Correct: \code{kubectl delete -f mongodb.yaml}
  (then remove from the list), or prune/GitOps machinery that
  reconciles deletions.
\end{enumerate}

\section*{Chapter 21 --- CI/CD Concepts}
\begin{enumerate}
  \item CI: main is always verified buildable --- the Build stage.
  CD: every verified change reaches production unaided --- the Deploy
  stage plus its rollout gate.
  \item Jenkins sits behind NAT with no public address --- GitHub
  cannot call in, so Jenkins polls. Same constraint: CI inside a
  corporate network that the cloud git host cannot reach.
  \item A rebuild is a different binary --- different dependency
  resolutions, timestamps, toolchain --- so tests certified something
  else. Chapter~1's env-var precedence lets one artifact take
  per-environment config at runtime.
  \item Tests skipped (re-enable, or gate elsewhere and know it);
  deploy not verified (rollout status/smoke tests); success defined
  as exit-codes not behavior (probe-backed verification, alerts).
\end{enumerate}

\section*{Chapter 22 --- Jenkins}
\begin{enumerate}
  \item Jobs, credentials, plugin state, build history. Image =
  reproducible binary+plugins; volume = irreplaceable state; together
  the server is cattle with a pet directory.
  \item Jenkins runs on the base image's JDK\,21; builds export
  \code{JAVA\_HOME=JDK25\_HOME} in the Jenkinsfile's environment
  block --- per-pipeline, not global.
  \item apply -k needs get/list + create/patch/update on every
  applied kind; rollout status needs get/list/watch on deployments
  and replicasets. Droppable without \code{set image}: nothing new
  --- set image's get+patch is already covered; the pods/log read
  rules exist only for the failure handler.
  \item Jib pushes to \code{localhost:5000}, and that name must mean
  the same machine to pusher and puller; host networking makes
  Jenkins' localhost the host's, which is what the image name
  promised k3s.
  \item ServiceAccount = identity; Role = capability list;
  RoleBinding = the join. Role-not-ClusterRole caps a compromised
  pipeline at one namespace --- blast radius as a type choice.
\end{enumerate}

\section*{Chapter 23 --- The Pipeline}
\begin{enumerate}
  \item SNAPSHOT is mutable --- a node could serve any historical
  build under the same name; the SHA pins the exact commit and makes
  ``what is running?'' answerable from \code{kubectl describe}.
  \item Build only listed modules plus their dependencies. A
  forgotten MODULES entry: the new service never gets an image; the
  deploy loop then fails (or deploys a stale image if one exists) ---
  the pipeline's one manual coupling.
  \item \code{kubectl rollout status --timeout=300s} --- it converts
  probe results (Chapter~17) into pipeline exit codes.
  \item All rollouts should proceed concurrently; waiting per-service
  serializes convergence that Kubernetes happily parallelizes ---
  same total work, longer wall clock.
  \item Changes: trigger (webhook), where builds run (ephemeral
  runners), secret mechanism (repo secrets). Unchanged: stage order,
  Maven commands, SHA tagging, the rollout verification.
\end{enumerate}

\section*{Chapter 24 --- Operations}
\begin{enumerate}
  \item get pods $\rightarrow$ describe pod $\rightarrow$ logs
  (\code{--previous} --- essential for CrashLoop, where the current
  container has no history yet) $\rightarrow$ exec and probe from
  inside.
  \item CrashLoop: the process exits --- application layer.
  \code{Running 0/1}: the process lives but readiness fails ---
  the probe's target (often a dependency) is the layer to inspect.
  \item Exec into the pod and try the connection from there --- the
  pod's DNS, network path, and injected env differ from your
  laptop's; your laptop proves only that Postgres is alive, not that
  the pod can reach it.
  \item k3s pod/service traffic traverses the host firewall;
  blocking \code{10.42.0.0/16} or \code{10.43.0.0/16} breaks DNS and
  pod-to-pod calls intermittently --- symptoms indistinguishable from
  flaky applications.
\end{enumerate}

\section*{Chapter 25 --- Kafka in the Cluster}
\begin{enumerate}
  \item A controller whose reconciliation loop targets custom resource
  kinds defined by CRDs. Strimzi added \code{Kafka}/\code{KafkaNodePool}
  (the cluster) and \code{KafkaTopic} (the contracts) --- plus user,
  connect and bridge kinds this project does not use.
  \item The operator bundle contains cluster-scoped objects (CRDs,
  ClusterRoles) that a namespaced Role cannot grant; \code{kafka.yaml}
  contains only namespaced custom resources, which the Role now
  permits. Role vs.\ ClusterRole is the line.
  \item Advertised listeners, the controller quorum voters, the
  per-listener security protocol map (also the cluster ID and per-broker
  Services). The operator computes them from the listener declaration
  and the node pool.
  \item The broker pod's actual status (\code{get pods}) and the
  operator log's reconciliation lines. The condition may be a stale
  write from a reconciliation that timed out during a long image pull;
  the next periodic loop corrects it.
  \item Reviewed partition counts, no day-one producer/topic race, and
  an inventory of event contracts as objects. One partition: total
  order per topic, sufficient while every consumer group has one
  instance; partitions can be added later but never removed.
\end{enumerate}

\section*{Chapter 26 --- The Frontend and the Single Host}
\begin{enumerate}
  \item Stage one (Node) compiles and bundles; stage two (nginx) serves.
  Only \code{dist/} is copied with \code{COPY --from=build}; Node,
  \code{node\_modules} and the compiler are discarded, leaving a
  $\sim$15\,MB image with a smaller attack surface.
  \item Traefik matches \code{/courses/7} against both Ingresses;
  \code{/} is the longest matching prefix (no \code{/courses} rule
  exists), so nginx receives it. No such file exists;
  \code{try\_files \$uri \$uri/ /index.html} falls back to
  \code{index.html}, the app boots and React Router renders the route
  from the URL.
  \item Rules for the same host are merged into one routing table and
  the longest matching path prefix wins: \code{/api/courses} matches
  \code{/api} (longer) over \code{/}; \code{/assets/x.js} matches only
  \code{/}.
  \item CORS applies only to cross-origin requests. In the cluster the
  page and the API share scheme, host and port (\code{http://ts.local}),
  so the browser never sends a preflight; on a laptop the page is
  \code{localhost:3000} and the API \code{localhost:8080} --- a
  different port is a different origin.
  \item The socket is owned by the host's \code{docker} group (GID 984
  on ts-server); the jenkins user inside the container is not in a
  group with that numeric ID, so writes are denied. \code{--group-add}
  attaches the host's GID at run time; computing it keeps the script
  correct on any host without hard-coding a machine-specific number.
\end{enumerate}

\section*{Chapter 27 --- The Road Ahead}
\begin{enumerate}
  \item Consumer-group lag: the distance between the newest offset in a
  partition and the group's committed offset --- growing lag means
  notification-service is falling behind or down.
  \item Push: the credential lives in CI (Jenkins' kubeconfig).
  Pull: it never leaves the cluster; CI only writes git. Pruning
  resolves Chapter~20's apply-never-deletes pitfall.
  \item Gateway routes change \code{lb://name} to
  \code{http://name:port}; Service endpoints replace registration,
  readiness probes replace heartbeats; delete discovery-server (and
  eureka client deps), free its resources.
  \item They already run as non-root with group-0 permissions, so
  OpenShift's arbitrary-UID SCC admits them unmodified; only
  Ingress $\rightarrow$ Route remains.
\end{enumerate}

\section*{Chapter 28 --- Deep Dive: JWT at Scale}
\begin{enumerate}
  \item No. The config-server now holds only the public key, which
  verifies but cannot sign. Forging requires the private PEM, which
  lives solely in auth-service's Secret --- that Secret (or the
  auth-service pod) is what would have to leak.
  \item Publish the new public key in the JWKS under a new \code{kid}
  while keeping the old; wait for verifier caches to refresh; switch
  signing to the new private key; after the longest token lifetime,
  remove the old key. In step~3 both generations of token are in
  flight at once; \code{kid} tells each verifier \emph{which} public
  key to apply, so old tokens keep verifying and the old key can later
  be retired with certainty rather than by trial.
  \item Users already holding an access token are unaffected until it
  expires (15 minutes), because verifiers cached the public key. Then
  \code{/refresh} fails and they are logged out; new logins fail from
  the first minute. A verifier restarted during the outage boots with
  an empty JWKS cache and, failing closed, rejects every token ---
  which is why the public key should also be available from
  configuration as a fallback.
  \item Staleness: the claim is frozen at issue time, so a course
  created after login is invisible until the token expires --- the
  token cannot know about data it does not own. Ownership: course
  membership is course-service's data; encoding it in an
  auth-service-issued token couples the two services' schemas through
  the token contract. (Size is the third reason, and the least
  important.)
  \item The header is forwarded, but no public endpoint reads
  \code{X-User-Id}: \code{/refresh} authenticates by refresh token,
  \code{/actuator} by nothing. So it cannot be exploited today, yet
  the invariant ``only the gateway writes these headers'' is
  unenforced. \code{default-filters: RemoveRequestHeader=X-User-Id}
  (and \code{X-User-Role}) in \code{api-gateway.yml} enforces it.
  \item Up to 15 minutes --- the access-token lifetime --- since the
  gateway verifies locally and only \code{/refresh} consults the
  database. A \code{jti} deny-list adds Redis (one lookup per request,
  Redis down means fail-open or fail-closed by your choice); the
  phantom-token pattern adds Redis on the critical path of every
  request with no fail-open option --- Redis down is a full outage.
\end{enumerate}

\section*{Chapter 29 --- Deep Dive: Configuration}
\begin{enumerate}
  \item The Secret wins: it arrives as the env var
  \code{SPRING\_DATASOURCE\_PASSWORD}, and OS environment outranks the
  config-server document. If the Secret object is missing,
  \code{envFrom} fails the pod; if only the key is missing, the pod
  starts, the served file's literal \code{root} takes effect, and the
  service connects as \code{root}, silently.
  \item Precedence chooses \emph{which document} defines
  \code{app.s3.secret-key}; the config-server's document does, and it
  sends the literal text \code{\${APP\_S3\_SECRET\_KEY:minioadmin}}
  unresolved. The client resolves the placeholder against its own
  Environment afterwards, so the env var wins because the text asked
  for it, not because of source order.
  \item Window one: between the two refreshes, auth-service signs with
  the new key while the gateway verifies with the old --- every login
  yields a token the gateway rejects. Window two: after both refresh,
  every token issued before the change is rejected until it would have
  expired anyway. Hot rotation needs overlapping keys selected by
  \code{kid}, which a single shared HS256 secret cannot express.
  \item Non-optional: the fetch fails, the JVM exits, kubelet restarts
  with growing back-off until the server returns. \code{optional:}:
  boots from packaged defaults pointing at localhost, passes its
  readiness probe, takes traffic, fails every request. \code{fail-fast}
  with retry: retries in-process with back-off for the configured
  window, then exits --- no wasted JVM starts, and no lying pod.
  \item Kustomize's generator appends a content hash to the ConfigMap
  name; a changed file is a new name, the Deployment referencing it is
  a new pod template, and Kubernetes rolls the pods. Config, code and
  rollback become one mechanism with history --- the server offered
  none of that without Bus or restarts.
  \item Remove it when every runtime is Kubernetes, no value must
  change without a pod restart, and secrets already live in the
  platform. Keep it when a non-Kubernetes runtime (VMs, Compose in
  production) must share the same configuration source.
\end{enumerate}

\section*{Chapter 30 --- Deep Dive: Service Discovery}
\begin{enumerate}
  \item Registration is immediate on web-server start (Spring forces
  it). Then the server's read-only response cache (30\,s), the
  gateway's registry fetch interval (30\,s), and the LoadBalancer cache
  TTL (35\,s): about 95\,s worst case, roughly half that typically.
  \item Lease expiry 90\,s, then the eviction sweep (up to 60\,s), then
  server cache 30\,s, client fetch 30\,s, LoadBalancer cache 35\,s ---
  up to 245\,s. A graceful shutdown sends DELETE and skips the first two
  (150\,s), leaving only the caches.
  \item Expected renewals per minute: $5 \times 2 = 10$; threshold
  $\lfloor 10 \times 0.85 \rfloor = 8$; the server stays healthy only
  while renewals exceed 8. One dead client gives exactly 8, so
  self-preservation engages and nothing is evicted.
  \code{eureka.server.enable-self-preservation: false} (or a lower
  \code{renewal-percent-threshold}) fixes it on a small fleet.
  \item Eureka's self-reported status gates the gateway path, because
  \code{lb://} resolves to a registered pod IP and never touches the
  Kubernetes Service where readiness is enforced. Routing to
  \code{http://course-service:8082} puts kube-proxy, and therefore
  readiness, back on the path.
  \item Compose activates no profile, so the base document applies:
  Eureka client enabled, \code{lb://} routes, and the Feign
  \code{url} resolves to an empty string, which Feign treats as ``use
  load-balanced resolution by name''. Under \code{k8s} the property is
  a Service URL and Feign bypasses the LoadBalancer entirely.
  \item The registry knows an instance only through what the instance
  says about itself, and learns of silence only after a lease it does
  not sweep immediately; a Service's endpoints are maintained by the
  kubelet, which actively probes the pod every few seconds and removes
  it the moment a probe fails. Self-report with a grace period versus
  external observation with none.
\end{enumerate}

\section*{Chapter 31 --- Deep Dive: The Gateway Under Load}
\begin{enumerate}
  \item Nothing has been sent yet: \code{chain.filter} returns a
  \code{Mono} describing the rest of the chain; it runs only when the
  framework subscribes, and the calling thread is already free. A
  servlet filter's \code{doFilter} blocks until the downstream call
  completes, so the thread is held for the whole request --- which is
  the model the event loop exists to avoid.
  \item Every request whose socket is owned by the stalled loop stops
  --- including routes that never touch Redis --- because loops are
  shared across routes, not per route. Reactor Netty runs
  \code{max(cores, 4)} loops; on a 2-core box that is four, so one
  blocked loop is a quarter of the gateway's capacity.
  \item Sixteen: the default per-host pool is \code{max(cores, 8)
  $\times$ 2}. With no response timeout, 16 hung requests hold all 16
  connections and the 17th waits 45\,s in the acquire queue. Setting
  \code{spring.cloud.gateway.httpclient.response-timeout} (or raising
  \code{pool.max-connections}) changes the number; the timeout is the
  one that fixes the cause.
  \item The filter wraps each call in a Resilience4j TimeLimiter whose
  default is 1\,s. The 1.5\,s upload is cut off, counted as a failure,
  and after enough of them the breaker opens against a healthy
  service. Set \code{resilience4j.timelimiter.instances.<name>.timeout-duration}
  at least as long as the response timeout.
  \item With two or more replicas an in-memory bucket gives each replica
  its own limit, so a user gets $N\times$ the quota. Redis holds one
  bucket per key, updated atomically by a Lua script. Trade-off: Redis
  is a dependency on every request; the filter's default when Redis is
  unreachable is to allow the request (fail-open), which must be a
  deliberate choice.
  \item Every request is non-simple, so in dev each API call is preceded
  by an \code{OPTIONS} preflight. \code{setMaxAge} lets the browser
  cache the preflight result (Chrome's default is 5\,s). In the cluster
  the frontend and \code{/api} share one origin behind Traefik, so no
  request is cross-origin and the CORS filter is never consulted.
\end{enumerate}

\section*{Chapter 32 --- Deep Dive: Synchronous Calls}
\begin{enumerate}
  \item (a) No code calls \code{AuthServiceClient} --- wire it into a
  service such as \code{UserLookupService}. (b) Feign's breaker
  integration is off --- set \code{circuitbreaker.enabled: true} under
  \code{spring.cloud.openfeign}.
  (c) No \code{CircuitBreakerFactory} bean exists --- add the
  \code{spring-cloud-starter-circuitbreaker-resilience4j} dependency.
  Separately, nothing references the \code{authService} instance ---
  annotate the calling method with \code{@CircuitBreaker(name =
  "authService")}.
  \item Nine: \code{minimum-number-of-calls} defaults to the window
  size, so the breaker does not evaluate until the tenth outcome. A
  30\,s call is a success because \code{slow-call-duration-threshold}
  defaults to 60\,s; only calls slower than that count as slow, and
  only if the slow-call rate threshold (default 100\,\%) is met.
  \item Browser (longest) $>$ gateway response timeout $>$ Feign read
  timeout $\times$ retry attempts (plus backoff). Each hop must give up
  before its caller does, or it works on abandoned requests; and the
  retry count multiplies the innermost timeout, so it must fit inside
  the gateway's budget --- 2\,s $\times$ 2 attempts under a 5\,s
  gateway timeout, not 3.
  \item A \code{NullPointerException} in every caller that forgets to
  check, surfacing only during an outage; and the inability to tell a
  404 (user genuinely missing) from a breaker rejection, so both get
  the same treatment. \code{degraded(Long id, Throwable cause)} ---
  the two-argument fallback --- receives the cause and can return a
  cached or placeholder value for infrastructure failures while letting
  a 4xx propagate.
  \item auth-service's \code{SecurityConfig} permits only public paths;
  \code{/users/\{id\}} needs an authenticated context that its
  \code{GatewayHeaderAuthFilter} builds from \code{X-User-Id}, which
  Feign does not send. Each 401 is an exception, so ten of them open
  the breaker against a healthy service. Fix: a Feign
  \code{RequestInterceptor} that forwards the identity headers, and
  \code{ignore-exceptions} for 4xx so client errors never count as
  dependency failures.
  \item $0.999^3 \approx 99.7\,\%$. Raise it by removing dependencies
  from the request path --- a local copy fed by events, or a cache in
  front of the call --- and by making the remaining ones optional with
  a typed fallback so their failure degrades the response instead of
  failing it.
\end{enumerate}

\section*{Chapter 33 --- Deep Dive: Data Across Services}
\begin{enumerate}
  \item Order A: the transaction commits, then the JVM dies before the
  producer buffer flushes --- user activated, no event, no student
  ever. Order B: the buffer flushes and the broker acks, then the
  commit fails --- course-service creates a student for a user that
  was rolled back. B produces the phantom row.
  \item The producer cannot fetch metadata, blocks for
  \code{max.block.ms} (60\,000\,ms), then throws; the exception leaves
  \code{register()}, the transaction rolls back, and the user is not
  created. With the outbox the request writes a user row and an outbox
  row in one commit and returns in milliseconds; the relay retries the
  send in the background.
  \item Two replicas run the same scheduler. Without the lock both
  select the same unpublished rows and both publish them --- a
  duplicate for every event, every second. \code{FOR UPDATE} makes the
  second replica wait; \code{SKIP LOCKED} makes it take the next rows
  instead, so they share the work without overlap.
  \item Ordering per key. If row~1 for user~42 fails and the loop
  \code{continue}s to row~2 for the same user, row~2 is published
  first; the next tick publishes row~1 behind it, and the consumer sees
  ``updated'' before ``activated''. Stopping the batch keeps the
  partition's order equal to the outbox's order.
  \item Nothing: \code{ts.user.activated} fires once, and no
  \code{ts.user.updated} exists, so \code{student.email} is frozen at
  activation. Changes: emit an outbox row with
  \code{eventType=UserUpdated} wherever the profile is saved; make
  \code{UserActivatedListener} upsert (find-or-create, then set
  fields); subscribe it to both topics or add a second listener.
  \item States: \code{ACTIVE} $\rightarrow$ \code{DELETING}
  $\rightarrow$ deleted, or back to \code{ACTIVE}. Events:
  \code{UserDeletionRequested} (auth, via outbox),
  \code{StudentDataDeleted} / \code{StudentDataDeletionFailed} (course,
  after removing student, enrollments, submissions, and MinIO objects).
  Compensation: on failure auth restores \code{ACTIVE} and alerts.
  \code{DELETING} is the semantic lock: it refuses logins and new
  enrollments while the saga runs, so no step operates on data another
  step is deleting.
\end{enumerate}

\section*{Chapter 34 --- Deep Dive: Kafka Delivery Guarantees}
\begin{enumerate}
  \item \code{enable.idempotence=true} (broker de-duplicates producer
  retries), \code{acks=all}, unbounded \code{retries} within
  \code{delivery.timeout.ms} --- all Kafka~3.x client defaults the
  project never set. On Compose there is one broker and
  \code{replicas: 1}, so ``all in-sync replicas'' is one disk; the
  setting only buys durability once \code{replication.factor=3} and
  \code{min.insync.replicas=2} exist.
  \item A batch is polled; the listener sends mails one by one on the
  poll thread; processing exceeds \code{max.poll.interval.ms}
  (300\,000\,ms); the broker evicts the consumer and reassigns the
  partition from the last committed offset; the new owner processes
  the same records while the old one finishes its loop.
  \item If the claim commits alone and the business write then fails,
  the retry finds the id already claimed, returns, and the offset
  advances --- the change is lost forever while the table says it
  happened. In one transaction, failure rolls back the claim too, so
  the retry can process it.
  \item A wait only helps transient faults: an SMTP relay coming back,
  a network blip. A \code{NullPointerException} or a deserialization
  error is deterministic; retrying it three times delays the same
  failure. It is excluded, so it goes straight to the DLT on the first
  attempt.
  \item Boot's \code{KafkaAutoConfiguration} backs off when a
  \code{ProducerFactory} bean exists, so \code{spring.kafka.producer.*}
  in \code{course-service.yml} is ignored by course-service. Fix:
  delete \code{KafkaProducerConfig} (Boot's factory reads the same
  serializers from YAML) --- or keep the bean and delete the YAML block.
  \item Retention is 168 hours; nine days is longer. The group's
  committed offset now points before the oldest retained record, so
  \code{auto-offset-reset: earliest} applies and it resumes from the
  oldest surviving record --- two days of events are gone. With a new
  group id there is no committed offset at all; \code{earliest} makes
  it process all seven retained days, re-sending every mail in them.
\end{enumerate}

\section*{Chapter 35 --- Deep Dive: Files at Scale}
\begin{enumerate}
  \item \code{endpointOverride}: the client uses \code{endpoint}
  (\code{minio:9000}), the presigner uses \code{public-endpoint}. If
  the presigner signed against the internal name, every download link
  would embed a host the browser cannot resolve --- and rewriting it
  afterwards breaks the signature.
  \item SigV4 signs a canonical request that includes the signed
  headers, and \code{Host} is always among them
  (\code{X-Amz-SignedHeaders=host}). A link signed for
  \code{minio:9000} and rewritten to another host presents a different
  canonical request; MinIO recomputes the signature, it does not match,
  and it returns \code{403 SignatureDoesNotMatch}.
  \item The transaction covers the Postgres insert only; the
  \code{putObject} is an HTTP call outside it. A failed \code{save}
  rolls back the row while the object remains in the bucket under a
  key no row records --- an orphan.
  \item Cheap: listing every attempt by one student for one assignment
  (one \code{ListObjectsV2} with the prefix). Expensive: finding all
  files by one student across assignments (no prefix matches; full
  scan or database query). The UUID makes every key unique, so no
  upload ever overwrites another; versioning would only protect against
  overwrites that cannot happen.
  \item Text submission succeeds (no storage call). Listing succeeds and
  returns links (presigning is local computation). Opening a link fails
  (MinIO must serve it). Restarting course-service fails: the
  \code{@PostConstruct} bucket check throws on connection refused.
  \item The key is built server-side from \code{assignmentId} and the
  caller's own \code{studentId}, and stored on a \code{PENDING\_FILE}
  row owned by that student; \code{completeUpload} checks that the row
  belongs to the caller and that a \code{HEAD} on \emph{that} row's key
  succeeds. A URL for another student's key cannot be attached to a
  different row because rows are never given a client-supplied key.
\end{enumerate}

\section*{Chapter 36 --- Deep Dive: The Dictionary on MongoDB}
\begin{enumerate}
  \item The \code{FETCH} stage's \code{filter} rejects them. The
  \code{createdByUserId} index only narrows to the user's slice; an
  unanchored, case-insensitive \code{\$regex} has no leading characters
  to seek by, so every fetched document is tested. Anchoring the pattern
  (\code{\^{}clou}) with a matching collation lets an index on
  \code{word} serve it as a range.
  \item All matches --- potentially all 3{,}000 --- come back from
  Mongo; \code{toResponse} then runs two \code{word\_links} queries per
  match (about 6{,}000); the page is built in Java by \code{subList}
  after sorting the full list.
  \item Multi-document transactions need a replica set; the project
  runs a standalone \code{mongod}. A crash after deleting the definition
  and before deleting its links leaves edges pointing at a missing word.
  \code{findByIdAndCreatedByUserId(...).orElse(null)} followed by the
  \code{filter(Objects::nonNull)} drops such nodes silently.
  \item \code{BsonMaximumSizeExceededException} (16\,MB limit). Six
  photos of about 2--3\,MB each, inflated by a third as base64, exceed
  16{,}793{,}600 bytes. Chapter~35's rule: bytes in object storage,
  keys in the database --- store the images in MinIO and keep their
  keys in \code{images}.
  \item A text index tokenizes and stems whole words; ``lou'' is not a
  token of ``cloud'', so there is no hit. Infix matching needs a regex
  scan or an n-gram/trigram index (Atlas Search, or \code{pg\_trgm} in
  Postgres).
  \item dictionary-service fails on every endpoint; every other service
  and topic is unaffected (no shared connection, no caller, no event).
  Adding a Mongo ping to the readiness probe would take the pod out of
  the gateway's endpoints, turning 500s into a clean 503 at the edge.
\end{enumerate}

\section*{Chapter 37 --- Deep Dive: Observability}
\begin{enumerate}
  \item Producer side:
  \code{spring.kafka.\allowbreak template.\allowbreak observation-enabled}
  (absent in every config document); consumer side:
  \code{spring.kafka.\allowbreak listener.\allowbreak observation-enabled}.
  Properties
  configure only the \emph{auto-configured} \code{KafkaTemplate};
  auth-service uses that one, so the property fixes it.
  Course-service builds its own template in \code{KafkaProducerConfig},
  which ignores the property and needs
  \code{setObservationEnabled(true)} in code.
  \item It continues the trace: Boot~3's default is to \emph{consume}
  W3C, B3 and B3-multi and to \emph{produce} W3C, decided by
  \code{management.tracing.propagation} defaults. The header leaving
  the gateway is \code{traceparent}, with the same trace id and a new
  parent span id.
  \item Exposure is a name list; the endpoint exists only when a
  registry that renders that format is on the classpath. No service
  depends on \code{micrometer-registry-prometheus}, so Actuator has
  nothing to expose under that name. Adding that one artifact (Boot
  BOM-managed) makes the existing exposure line real.
  \item Latency: unchanged --- spans are queued and shipped by a
  background thread. Heap: bounded --- the queue holds 10{,}000 spans
  and drops on overflow with a counted warning. Spans: those produced
  during the hour are lost; there is no replay. With a synchronous
  reporter, \emph{latency} would change: every request would block on
  a failing HTTP call to Zipkin until its timeout.
  \item One percent --- head-based sampling decides at the root before
  the outcome is known, so failures are kept at the same rate as
  successes. Raising it to 1.0 for errors only requires tail-based
  sampling, which needs an OpenTelemetry Collector buffering complete
  traces before the store; the direct Brave-to-Zipkin path cannot do
  it.
  \item An e-mail is an unbounded tag value; each distinct value
  creates a new time series in the metrics registry and a new indexed
  value in the trace store. Memory grows with the user base until
  \code{/actuator/metrics} crawls and the pod is \code{OOMKilled}. Rule:
  tags carry \emph{templates and enums} (route template, status,
  outcome); identifiers go into span annotations or log lines, which
  are not indexed by value.
\end{enumerate}

\section*{Chapter 38 --- Deep Dive: The Browser Meets the Gateway}
\begin{enumerate}
  \item Key \code{auth-storage} in \code{localStorage}, written by
  Zustand's \code{persist}. One line:
  \code{JSON.parse(localStorage.\allowbreak getItem('auth-storage')).\allowbreak
  state.refreshToken}.
  With it, the script (or whoever it sends it to) can call
  \code{/api/auth/refresh} for seven days and mint fresh access tokens
  at will, from any machine.
  \item A 401 from \code{/refresh} through \code{apiClient} would enter
  the response interceptor; \code{\_retry} is unset on that new
  request, so it would read the refresh token and POST \code{/refresh}
  again through \code{apiClient}; that 401 enters the interceptor
  again --- an unbounded chain of refresh calls, each spawning the
  next, until the tab is closed. Bare \code{axios} has no interceptor,
  so its 401 reaches the \code{catch} and logs out.
  \item All four get 401 and each sets its own \code{\_retry}; four
  POSTs to \code{/refresh} race. The first succeeds and rotates the
  token; the second presents the now-invalid old token and gets 401
  from auth-service; that lands in the \code{catch} --- step 4,
  \code{logout()} and a hard redirect --- while the first request's
  replay may even succeed. Single-flight: one promise, started by the
  first 401, awaited by the rest; one POST, one rotation, four
  replays with the same new token.
  \item (a) A vulnerable or forgotten endpoint elsewhere under
  \code{/api} never receives the cookie, so a CSRF or log-leak bug
  there cannot touch it; (b) the cookie is not attached to ordinary
  API calls, so it never appears in gateway access logs, proxies or
  browser traffic for \code{/api/courses}. The browser matches
  \code{Path} against the URL \emph{it} requests, which is the
  gateway's \code{/api/auth/refresh}; auth-service's stripped
  \code{/refresh} is never a browser-visible path.
  \item In neither: Vite's dev proxy and Traefik in the cluster both
  keep page and API on one origin, so no request is cross-origin and
  no preflight occurs. Serving the built bundle from a different
  origin --- a CDN, or a static host on another port without the
  proxy --- would make every API call cross-origin and the config
  decisive.
  \item CORS is enforced by the browser after the response arrives:
  the gateway processed the request and returned 200, then the browser
  refused to hand the body to the script because the response lacked
  an acceptable \code{Access-Control-Allow-Origin}. Compare: the
  request's \code{Origin} header versus the allowed-origins list (exact
  string, scheme and port); \code{Access-Control-Allow-Credentials}
  versus whether the request was credentialed; and for the preflight,
  \code{Access-Control-Allow-Headers} versus the headers the request
  actually sends (\code{Authorization}, \code{Content-Type}).
\end{enumerate}
